\editionPage{indices-defis}{Indications progressives}{Défis D01 à D10}{Lire une piste, puis reprendre sa recherche.}\fontsize{10}{13}\selectfont
\par\addvspace{3mm}\noindent\begin{minipage}{\linewidth}
{\otSemi\color{otNavy}Défi D01}\phantomsection\label{indice-1}\par\vspace{1mm}
Piste 1 : séparer $s$ et $u$. Piste 2 : utiliser les racines de 16 et 36. Piste 3 : examiner la référence après les calculs.
\end{minipage}\par

\par\addvspace{3mm}\noindent\begin{minipage}{\linewidth}
{\otSemi\color{otNavy}Défi D02}\phantomsection\label{indice-2}\par\vspace{1mm}
Piste 1 : soustraire la tare. Piste 2 : identifier le dénominateur de chaque quotient. Piste 3 : comparer des grandeurs de même nature.
\end{minipage}\par

\par\addvspace{3mm}\noindent\begin{minipage}{\linewidth}
{\otSemi\color{otNavy}Défi D03}\phantomsection\label{indice-3}\par\vspace{1mm}
Piste 1 : conserver la masse de soluté. Piste 2 : combiner des prélèvements. Piste 3 : contrôler le facteur de dilution.
\end{minipage}\par

\par\addvspace{3mm}\noindent\begin{minipage}{\linewidth}
{\otSemi\color{otNavy}Défi D04}\phantomsection\label{indice-4}\par\vspace{1mm}
Piste 1 : $A=Z+N_{\mathrm{neutrons}}$. Piste 2 : $q=(Z-N_e)e$. Piste 3 : sommer les charges.
\end{minipage}\par

\par\addvspace{3mm}\noindent\begin{minipage}{\linewidth}
{\otSemi\color{otNavy}Défi D05}\phantomsection\label{indice-5}\par\vspace{1mm}
Piste 1 : remplir les deux premières couches. Piste 2 : placer sept électrons dans la troisième. Piste 3 : comparer les noyaux.
\end{minipage}\par

\par\addvspace{3mm}\noindent\begin{minipage}{\linewidth}
{\otSemi\color{otNavy}Défi D06}\phantomsection\label{indice-6}\par\vspace{1mm}
Piste 1 : utiliser tous les indices de la formule. Piste 2 : concentration fois volume. Piste 3 : passer aux kilogrammes avant de compter.
\end{minipage}\par

\par\addvspace{3mm}\noindent\begin{minipage}{\linewidth}
{\otSemi\color{otNavy}Défi D07}\phantomsection\label{indice-7}\par\vspace{1mm}
Piste 1 : toute la masse change d’état. Piste 2 : diviser l’énergie par l’apport par seconde. Piste 3 : examiner la fin du palier.
\end{minipage}\par

\par\addvspace{3mm}\noindent\begin{minipage}{\linewidth}
{\otSemi\color{otNavy}Défi D08}\phantomsection\label{indice-8}\par\vspace{1mm}
Piste 1 : lire les rapports des coefficients. Piste 2 : calculer la consommation d’oxygène. Piste 3 : faire un bilan final.
\end{minipage}\par

\par\addvspace{3mm}\noindent\begin{minipage}{\linewidth}
{\otSemi\color{otNavy}Défi D09}\phantomsection\label{indice-9}\par\vspace{1mm}
Piste 1 : séparer identité et pureté. Piste 2 : attribuer chaque tache. Piste 3 : proposer un contrôle comparable.
\end{minipage}\par

\par\addvspace{3mm}\noindent\begin{minipage}{\linewidth}
{\otSemi\color{otNavy}Défi D10}\phantomsection\label{indice-10}\par\vspace{1mm}
Piste 1 : deux bilans indépendants. Piste 2 : comparer les tailles des noyaux. Piste 3 : énergie totale divisée par énergie par événement.
\end{minipage}\par

\editionPage{indices-suite}{Indications progressives}{Défis D11 à D20}{Lire une piste, puis reprendre sa recherche.}\fontsize{10}{13}\selectfont
\par\addvspace{3mm}\noindent\begin{minipage}{\linewidth}
{\otSemi\color{otNavy}Défi D11}\phantomsection\label{indice-11}\par\vspace{1mm}
Piste 1 : calculer chaque déplacement. Piste 2 : employer une échelle de vitesse distincte. Piste 3 : discuter ce qui est observé entre deux mesures.
\end{minipage}\par

\par\addvspace{3mm}\noindent\begin{minipage}{\linewidth}
{\otSemi\color{otNavy}Défi D12}\phantomsection\label{indice-12}\par\vspace{1mm}
Piste 1 : utiliser la dépendance en $1/d^2$. Piste 2 : partir du poids au sol. Piste 3 : distinguer gravitation et contact avec un support.
\end{minipage}\par

\par\addvspace{3mm}\noindent\begin{minipage}{\linewidth}
{\otSemi\color{otNavy}Défi D13}\phantomsection\label{indice-13}\par\vspace{1mm}
Piste 1 : séparer force et vitesse. Piste 2 : observer autour du sommet. Piste 3 : préciser le protocole de pointage.
\end{minipage}\par

\par\addvspace{3mm}\noindent\begin{minipage}{\linewidth}
{\otSemi\color{otNavy}Défi D14}\phantomsection\label{indice-14}\par\vspace{1mm}
Piste 1 : le son parcourt deux longueurs. Piste 2 : comparer sans arrondir trop tôt. Piste 3 : discuter l’origine du signal.
\end{minipage}\par

\par\addvspace{3mm}\noindent\begin{minipage}{\linewidth}
{\otSemi\color{otNavy}Défi D15}\phantomsection\label{indice-15}\par\vspace{1mm}
Piste 1 : diviser la durée totale par quatre. Piste 2 : distinguer oscillation entière et demi-oscillation. Piste 3 : séparer hauteur, timbre et intensité.
\end{minipage}\par

\par\addvspace{3mm}\noindent\begin{minipage}{\linewidth}
{\otSemi\color{otNavy}Défi D16}\phantomsection\label{indice-16}\par\vspace{1mm}
Piste 1 : séparer fond et raies. Piste 2 : compter les radiations. Piste 3 : chercher une référence physique.
\end{minipage}\par

\par\addvspace{3mm}\noindent\begin{minipage}{\linewidth}
{\otSemi\color{otNavy}Défi D17}\phantomsection\label{indice-17}\par\vspace{1mm}
Piste 1 : appliquer la loi à chaque indice. Piste 2 : comparer les angles à la normale. Piste 3 : distinguer angle et séparation sur l’écran.
\end{minipage}\par

\par\addvspace{3mm}\noindent\begin{minipage}{\linewidth}
{\otSemi\color{otNavy}Défi D18}\phantomsection\label{indice-18}\par\vspace{1mm}
Piste 1 : tracer le rayon central et celui qui arrive parallèlement à l’axe. Piste 2 : lire les deux échelles. Piste 3 : comparer au cas du chapitre.
\end{minipage}\par

\par\addvspace{3mm}\noindent\begin{minipage}{\linewidth}
{\otSemi\color{otNavy}Défi D19}\phantomsection\label{indice-19}\par\vspace{1mm}
Piste 1 : même tension dans les branches. Piste 2 : appliquer Ohm puis la loi des nœuds. Piste 3 : vérifier le branchement du multimètre.
\end{minipage}\par

\par\addvspace{3mm}\noindent\begin{minipage}{\linewidth}
{\otSemi\color{otNavy}Défi D20}\phantomsection\label{indice-20}\par\vspace{1mm}
Piste 1 : convertir l’entier en tension. Piste 2 : appliquer le pont puis l’étalonnage. Piste 3 : contrôler les domaines de chaque relation.
\end{minipage}\par

